\documentclass[a4paper,12pt,chapterprefix=false,twoside,openright]{scrbook}
\usepackage[utf8]{inputenc}
\usepackage{fontspec}
\usepackage[headings]{fullpage}
\usepackage{helvet}


%Grafikeinbindung
\usepackage{graphics,xcolor}
%\usepackage[pdftex]{graphicx,xcolor}
\graphicspath{{./src/pics/}}
\usepackage{pgfplots}
\usepackage{tikz} 
\usetikzlibrary{plotmarks,shapes,arrows.meta,decorations.markings,patterns,calc}

%Mathematikfunktionen
\usepackage{amsmath,amsfonts,amssymb}
\usepackage{mathtools}

%Tabellen
\usepackage{tabularx}
\usepackage{booktabs}

%Setup a Times New Roman Environment for the Title Page
\newenvironment{ptm}{\fontfamily{ptm}\fontsize{20}{24}\selectfont}{\par}
%Used on title page
\newcolumntype{K}[1]{>{\centering\arraybackslash}b{#1}}

\makeindex

%Headers
\usepackage{scrlayer-scrpage}
\pagestyle{headings}

% Include paper
\usepackage{pdfpages}
\usepackage{xb}
  
%==================================================================
 %==================================================================
 % Metadaten der Arbeit
 %
 % Titel
 \newcommand{\ifrtitle}{A Master thesis template with conference paper}
 % Author
  \newcommand{\ifrauthor}{Firstname Lastname}
 % Typ der Arbeit
   \newcommand{\ifrtype}{Master Thesis}
 % Betreuer
  \newcommand{\ifrsupervisor}{Dr. Torbjørn Cunis}
 %
 %=================================================================== 
 %==================================================================
  
  
\usepackage[pdfauthor={\ifrauthor},
          pdftitle={\ifrtitle},
          pdfsubject={\ifrtype},
          pdfkeywords={},
          pdfproducer={pdflatex},
          pdfcreator={Latex}]{hyperref}
  

%Beispieltext
\usepackage{blindtext}

%%Für Automatische Nomenklatur folgenden Block einkommentierten: UND ZEILE 113
%Für Hinweise bitte ./src/parts/frontmatter/readme_nomenclature.txt lesen.
%==================================================================
%Einheiten
\usepackage{units}
%%Nomenklauturpaket
\usepackage[intoc]{nomencliFR_A4}
\renewcommand{\nomname}{}
%Definition der Einheitenspalte
\newcommand{\nomenunit}[1]{\hskip-0em\parbox{4em}{$\left[ #1 \right]$}}

\makenomenclature
%%==================================================================


\begin{document}

%Frontmatter
%=============================================================================
\frontmatter

%Title
%********************************
% Titelseite
%********************************
\begin{titlepage}

\vspace*{2cm}
\begin{center}
\begin{ptm}
\fboxrule1.0mm
\fbox{
\begin{minipage}{\textwidth}
\begin{center}
\vspace{0.5cm}
 \textbf{\ifrtitle}
\vspace{0.5cm}
\end{center}
\end{minipage}
}

\vspace{4cm}
\textbf{\ifrauthor}\\
\vspace{1cm}
\textbf{\ifrtype}\\
\vspace{0.5cm}
\textbf{\the\year}
\\

\end{ptm}
\end{center}
\vfill
\begin{tabular}{K{2.7cm}K{13.0cm}}
	\begin{minipage}{\columnwidth}
	 \includegraphics[scale=0.4, trim= -0.0cm 0.05cm 0cm 0cm] {./src/pics/logo/ifrlogo.pdf}
	 \vspace{0.1cm}
	 \end{minipage}
	  & 
	\begin{minipage}{\columnwidth}
	\vfill
 \sffamily \large Universität Stuttgart \\ \textbf{Institut für Flugmechanik und Flugregelung}\vspace{0.1cm}
	 \end{minipage}\\
	 
	 \hline
 \end{tabular}

\end{titlepage}
\cleardoublepage

%Erklärung
\addchap*{Erklärung}
\thispagestyle{empty}
Hiermit versichere ich, dass ich diese \ifrtype~selbständig mit Unterstützung der Betreuer angefertigt und keine anderen als die angegebenen Quellen und Hilfsmittel verwendet habe.\\
\ \\
Die Arbeit oder wesentliche Bestandteile davon sind weder an dieser noch an einer anderen Bildungseinrichtung bereits zur Erlangung eines Abschlusses eingereicht worden.\\
\ \\
Ich erkläre weiterhin, bei der Erstellung der Arbeit die einschlägigen Bestimmungen zum Urheberschutz fremder Beiträge entsprechend den Regeln guter wissenschaftlicher Praxis\footnote{Nachzulesen in den DFG-Empfehlungen zur \glqq Sicherung guter wissenschaftlicher Praxis\grqq~bzw. in der Satzung der Universität Stuttgart zur \grqq Sicherung der Integrität wissenschaftlicher Praxis und zum Umgang mit Fehlverhalten in der Wissenschaft\grqq~} eingehalten zu haben. Soweit meine Arbeit fremde Beiträge (z.B. in der Form von Bilder, Zeichnungen, Textpassagen etc.) enthält, habe ich diese Beiträge als solche gekennzeichnet (Zitat, Quellenangabe) und eventuell erforderlich gewordene Zustimmungen der Urheber zur Nutzung dieser Beiträge in meiner Arbeit eingeholt. Mir ist bekannt, dass ich im Fall einer schuldhaften Verletzung dieser Pflichten die daraus entstehenden Konsequenzen zu tragen habe.\\[2cm]
\parbox{8cm}{
\hrule
\strut \centering\footnotesize Ort, Datum}
\hfill
\parbox{6cm}{
\hrule
\strut \centering\footnotesize \ifrauthor}
\vfill 
\textbf{Betreuer: \ifrsupervisor}
\cleardoublepage

%Abstract
~\vfill
\section*{Kurzfassung}
% Hier Text einfügen
\blindtext

\section*{Abstract}

% Hier Text einfügen
\blindtext

\vfill~
\cleardoublepage

%Inhaltsverzeichnis
\thispagestyle{empty}
\tableofcontents
\cleardoublepage

%Für Automatische Nomenklatur folgenden Zeile einkommentierten:
\chapter*{Nomenclature}
\addcontentsline{toc}{chapter}{Nomenclature}

%Greek
\printnomenclature{G}
%Latin
\printnomenclature{L}
%Subscripts
\subsubsection*{Indices}
\printnomenclature{S}
%Acronyms
\subsubsection*{Acronyms}
\printnomenclature{A}

\cleardoublepage

%%Samples for including Symbols etc. in the notation
\nomenclature[A]{KOS}{Koordinatensystem}

\nomenclature[L]{$_y\boldsymbol{T}_x$}{\nomenunit{-}Transformationsmatrix von System x nach System y}
\nomenclature[L]{$\boldsymbol{r}$}{\nomenunit{-}Vektor vom Lasthaken zum Schwerpunkt der Last}

\nomenclature[G]{$\Omega$}{\nomenunit{\unitfrac{rad}{s}} Drehgeschwindigkeit des Rotors}
\nomenclature[G]{$\Phi$,$\Theta$,$\Psi$}{\nomenunit{\unit{rad}}Rolll-, Nick- und Gierwinkel des Hubschraubers}

\nomenclature[S]{$_a(\cdot)$}{Aufhängepunktfestes Koordinatensystem}
\nomenclature[S]{$_g(\cdot)$}{Geodätisches Koordinatensystem}
\nomenclature[S]{$(\cdot)_L$}{Die Last, bzw. Außenlast betreffend}


%Mainmatter
%=============================================================================
\mainmatter

% !TeX root = ../../../Main.tex

\chapter{Introduction}

% !TeX root = ../../../Main.tex
\chapter{Background}

% !TeX root = ../../../Main.tex

%\chapter{Kapitel 3}
\chapter{Unknown}
% !TeX root = ../../../Main.tex
\chapter{Novel Algorithm}

%% !TeX root = ../../../Main.tex
%\chapter{Kapitel 5}
%% !TeX root = ../../../Main.tex
\chapter{Kapitel 6}
% !TeX root = ../../../Main.tex
\chapter{Conclusions}
\blindtext

%Backmatter
%=============================================================================
\backmatter

%Abbildungs- und Tabellenverzeichnis
\addcontentsline{toc}{chapter}{Abbildungsverzeichnis}
\listoffigures
\cleardoublepage
\addcontentsline{toc}{chapter}{Tabellenverzeichnis}
\listoftables

% Bibliographie
\addcontentsline{toc}{chapter}{Literaturverzeichnis}
\bibliographystyle{abbrvdin}
\bibliography{./src/bib/bibfile,./src/bib/IEEEexample}

%Appendix
%=============================================================================
\appendix

% first appendix
% !TeX root = ../../../Main.tex

\chapter{Appendix}
\blindtext


\end{document}